\section{Introduction}
\label{sec:introduction}

Probabilistic circuits (PCs) are density models that are expressive and tractable at the same time \citep{pc_intro}.
Because their structure is constrained, a PC answers marginal and conditional queries exactly and in time linear in its size.
For data with missing features this is exactly what one wants: a PC trained on incomplete data integrates the missing features out instead of guessing them.

When the task is classification, however, PCs are rarely combined with the classifier that is most often used on top of features, a logistic regression (LR).
The reason is the sigmoid.
Given a partially observed input $\xo$, the principled prediction of an LR is its expected prediction under the PC,
\begin{equation}
  p(y{=}1 \mid \xo) = \E_{p(\xm \mid \xo)} \left[ \sigma(w^T x) \right],
  \label{eq:goal}
\end{equation}
and the sigmoid does not factorise over the features, so the expectation cannot be pushed into the leaves of the circuit.
Even for a naive Bayes distribution over binary features, computing \eqref{eq:goal} is NP-hard \citep{khosravi2019expect}.
The tractability of the PC is lost as soon as the sigmoid is put on top.

The usual way to classify with a PC is therefore to avoid the LR altogether: train one joint model $p(x, y)$ and compute the conditional $p(y \mid x)$ by Bayes' rule.
This keeps everything tractable, but the parameters of the joint model are fitted to explain $x$, not to separate the classes, and in practice the resulting classifier is often worse than a discriminative one.

% TODO: our approach (CFI: the sigmoid as the CDF of the logistic distribution, Gil-Pelaez inversion of the
% characteristic function of z - L, the trapezoid rule), PCLR (PC and LR as one joint model trained jointly on
% incomplete data), and the contributions.

\section{A Logistic Regression on Top of a Probabilistic Circuit}
\label{sec:setup}

\paragraph{Probabilistic circuits.}
A PC is a computational graph of sum nodes, product nodes and leaves, where every leaf is a simple distribution over a single feature, for example a Gaussian \citep{pc_intro, spn}.
Every node has a scope, the set of features it is defined over.
A PC is \emph{smooth} if the children of every sum node have the same scope, and \emph{decomposable} if the children of every product node have disjoint scopes.
These two properties make expectations of functions that factorise over the features tractable: for $f(x) = \prod_d f_d(x_d)$, the expectation $\E[f(x)]$ is computed by evaluating the circuit once with every leaf replaced by $\E_{\text{leaf}}[f_d(x_d)]$.
With $f_d = 1$ for the missing features this gives the marginal $p(\xo)$, and with it the conditional $p(\xm \mid \xo)$.
A Gaussian mixture model (GMM) with diagonal covariances is the simplest such circuit, a single sum node over products of Gaussian leaves.

\paragraph{The joint model.}
We split an input $x \in \R^d$ into its observed features $\xo$ and its missing features $\xm$, and the weights $w$ of an LR accordingly into $w_o$ and $w_m$, so that the score is
\[
  z = w^T x = w_o^T \xo + w_m^T \xm.
\]
The joint model consists of a PC $p_\theta(x)$ over the features and an LR $p_w(y{=}1 \mid x) = \sigma(w^T x)$, with separate parameters $\theta$ and $w$.
For a complete input it predicts like the LR.
For an incomplete input it predicts with the expected prediction \eqref{eq:goal}, the average of the LR over all values of the missing features that the PC considers plausible.
The obvious way to compute it is Monte Carlo (MC): draw $\xm \sim p_\theta(\xm \mid \xo)$ and average $\sigma(w^T x)$.
This needs many samples for a precise answer, and it is random.
Conditional mean imputation, $\sigma(w_o^T \xo + w_m^T \E[\xm \mid \xo])$, is cheap but biased, since the sigmoid is not linear \citep{khosravi2019expect}.

% TODO: CFI, from the sigmoid to the characteristic function (the thesis methods chapter, Sections "Core Idea"
% and "Calculating the Integral").

\section{The Characteristic Function of a Circuit}
\label{sec:pc-cf}

The CF is the expectation of $e^{i s^T x} = \prod_d e^{i s_d x_d}$, a function that factorises over the features, so a smooth and decomposable PC computes it in one pass.
For the conditional CF we write
\[
  \phi_{\xm \mid \xo}(s) = \frac{1}{p(\xo)} \int e^{i s^T \xm}\, p(\xo, \xm)\, \mathrm{d}\xm.
\]
The integral is one pass through the circuit in which every leaf of an observed feature returns its density at $x_d$, as in a likelihood evaluation, and every leaf of a missing feature returns its own CF at $s_d$, for a Gaussian leaf $\exp(i s_d \mu - s_d^2 \sigma^2 / 2)$.
Dividing by the marginal $p(\xo)$, a pass with the missing leaves set to one, gives the conditional CF.
The values are complex, and we propagate them in log space, where products become sums and sums a complex log-sum-exp.
For a diagonal GMM this pass reduces to the closed form
\[
  \phi_{\xm \mid \xo}(s) = \sum_{k=1}^{K} \tilde{\pi}_k(\xo)\, e^{i s^T \mu_k^m - \frac{1}{2} s^T \Sigma_k^m s},
\]
with the responsibilities $\tilde{\pi}_k(\xo) \propto \pi_k\, \mathcal{N}(\xo; \mu_k^o, \Sigma_k^o)$ and the means and covariances $\mu_k^m, \Sigma_k^m$ of the components restricted to the missing features.
The responsibilities depend only on $\xo$, so they are computed once per input and reused for every frequency.

% TODO: the mean and variance of z for the step of CFI, by the same kind of pass (thesis methods chapter,
% "Probabilistic Circuits").

\section{Training the Joint Model on Incomplete Data}
\label{sec:training}

% TODO: PCLR, the joint model trained on incomplete data (thesis methods chapter, "PCLR: Training the Density and
% the Classifier Jointly").

\section{Experiments}
\label{sec:experiments}

% TODO: the estimation experiment (code/estimate.py) and the training experiment on make_classification
% (code/train.py, comparison and learning curve). The results so far are from a GMM that collapses to one
% component from 64 features on; they have to be rerun with the fixed GMM.

\subsection{Real Data}
\label{sec:results-real}

Six real datasets with continuous features and a binary label from OpenML: MAGIC (19{,}020 points, 10 features), Phoneme (5{,}404, 5), Spambase (4{,}601, 57), Banknote (1{,}372, 4), QSAR biodegradation (1{,}055, 41) and Pima diabetes (768, 8).
Every dataset is split 80/20 and standardised on the training points, and every entry is missing independently with probability 10, 30 or 50\,\%, in training and test.
Every model is trained on the incomplete training points and predicts the incomplete test points with its own rule: mean imputation, conditional mean imputation under a GMM with $K = 4$ components, MICE \citep{mice}, the generative classifier (a GMM per class and Bayes' rule), and PCLR.
The reference \emph{No missing} is the LR trained and tested on the complete points.
\Cref{tab:real-log-loss-10,tab:real-log-loss-30,tab:real-log-loss-50,tab:real-accuracy-10,tab:real-accuracy-30,tab:real-accuracy-50} give the mean and standard deviation over 5 seeds, each a new split and new missing entries; the best model per row (not counting the reference) is bold.

\begin{table*}[t]
  \centering
  \caption{Log-loss on the incomplete test points, 10\,\% missing.}
  \label{tab:real-log-loss-10}
  \resizebox{0.9\textwidth}{!}{\begin{tabular}{lcccccc}
\toprule
Dataset & No missing & Mean imp. & Cond.\ mean imp. & MICE & Generative & PCLR \\
\midrule
magic & 0.483{\scriptsize$\pm$0.005} & 0.509{\scriptsize$\pm$0.004} & 0.498{\scriptsize$\pm$0.003} & 0.500{\scriptsize$\pm$0.003} & 0.512{\scriptsize$\pm$0.005} & \textbf{0.497{\scriptsize$\pm$0.003}} \\
phoneme & 0.584{\scriptsize$\pm$0.007} & 0.594{\scriptsize$\pm$0.008} & 0.590{\scriptsize$\pm$0.008} & 0.594{\scriptsize$\pm$0.008} & \textbf{0.446{\scriptsize$\pm$0.008}} & 0.589{\scriptsize$\pm$0.009} \\
spambase & 0.258{\scriptsize$\pm$0.018} & \textbf{0.266{\scriptsize$\pm$0.024}} & 0.268{\scriptsize$\pm$0.022} & 0.272{\scriptsize$\pm$0.026} & 0.402{\scriptsize$\pm$0.028} & 0.271{\scriptsize$\pm$0.024} \\
banknote & 0.071{\scriptsize$\pm$0.009} & 0.180{\scriptsize$\pm$0.014} & \textbf{0.135{\scriptsize$\pm$0.011}} & 0.136{\scriptsize$\pm$0.005} & 0.334{\scriptsize$\pm$0.071} & 0.143{\scriptsize$\pm$0.013} \\
qsar & 0.474{\scriptsize$\pm$0.076} & 0.528{\scriptsize$\pm$0.067} & 0.516{\scriptsize$\pm$0.065} & \textbf{0.492{\scriptsize$\pm$0.065}} & 0.570{\scriptsize$\pm$0.029} & 0.528{\scriptsize$\pm$0.085} \\
pima & 0.563{\scriptsize$\pm$0.058} & 0.572{\scriptsize$\pm$0.054} & 0.567{\scriptsize$\pm$0.052} & 0.571{\scriptsize$\pm$0.059} & 0.590{\scriptsize$\pm$0.009} & \textbf{0.566{\scriptsize$\pm$0.052}} \\
\bottomrule
\end{tabular}
}
\end{table*}

\begin{table*}[t]
  \centering
  \caption{Log-loss on the incomplete test points, 30\,\% missing.}
  \label{tab:real-log-loss-30}
  \resizebox{0.9\textwidth}{!}{\begin{tabular}{lcccccc}
\toprule
Dataset & No missing & Mean imp. & Cond.\ mean imp. & MICE & Generative & PCLR \\
\midrule
magic & 0.483{\scriptsize$\pm$0.005} & 0.551{\scriptsize$\pm$0.006} & 0.524{\scriptsize$\pm$0.004} & 0.541{\scriptsize$\pm$0.005} & 0.523{\scriptsize$\pm$0.004} & \textbf{0.519{\scriptsize$\pm$0.004}} \\
phoneme & 0.584{\scriptsize$\pm$0.007} & 0.613{\scriptsize$\pm$0.008} & 0.605{\scriptsize$\pm$0.009} & 0.612{\scriptsize$\pm$0.009} & \textbf{0.474{\scriptsize$\pm$0.013}} & 0.599{\scriptsize$\pm$0.010} \\
spambase & 0.258{\scriptsize$\pm$0.018} & \textbf{0.307{\scriptsize$\pm$0.018}} & 0.314{\scriptsize$\pm$0.014} & 0.333{\scriptsize$\pm$0.030} & 0.409{\scriptsize$\pm$0.026} & 0.309{\scriptsize$\pm$0.015} \\
banknote & 0.071{\scriptsize$\pm$0.009} & 0.336{\scriptsize$\pm$0.024} & \textbf{0.280{\scriptsize$\pm$0.031}} & 0.289{\scriptsize$\pm$0.018} & 0.352{\scriptsize$\pm$0.025} & 0.285{\scriptsize$\pm$0.022} \\
qsar & 0.474{\scriptsize$\pm$0.076} & 0.563{\scriptsize$\pm$0.047} & 0.532{\scriptsize$\pm$0.031} & \textbf{0.499{\scriptsize$\pm$0.033}} & 0.537{\scriptsize$\pm$0.023} & 0.528{\scriptsize$\pm$0.022} \\
pima & 0.563{\scriptsize$\pm$0.058} & 0.580{\scriptsize$\pm$0.037} & 0.578{\scriptsize$\pm$0.041} & 0.581{\scriptsize$\pm$0.045} & 0.596{\scriptsize$\pm$0.021} & \textbf{0.575{\scriptsize$\pm$0.042}} \\
\bottomrule
\end{tabular}
}
\end{table*}

\begin{table*}[t]
  \centering
  \caption{Log-loss on the incomplete test points, 50\,\% missing.}
  \label{tab:real-log-loss-50}
  \resizebox{0.9\textwidth}{!}{\begin{tabular}{lcccccc}
\toprule
Dataset & No missing & Mean imp. & Cond.\ mean imp. & MICE & Generative & PCLR \\
\midrule
magic & 0.483{\scriptsize$\pm$0.005} & 0.592{\scriptsize$\pm$0.008} & 0.556{\scriptsize$\pm$0.005} & 0.577{\scriptsize$\pm$0.004} & \textbf{0.541{\scriptsize$\pm$0.003}} & 0.544{\scriptsize$\pm$0.007} \\
phoneme & 0.584{\scriptsize$\pm$0.007} & 0.631{\scriptsize$\pm$0.009} & 0.617{\scriptsize$\pm$0.009} & 0.634{\scriptsize$\pm$0.010} & \textbf{0.499{\scriptsize$\pm$0.009}} & 0.603{\scriptsize$\pm$0.010} \\
spambase & 0.258{\scriptsize$\pm$0.018} & \textbf{0.352{\scriptsize$\pm$0.022}} & 0.378{\scriptsize$\pm$0.012} & 0.367{\scriptsize$\pm$0.024} & 0.447{\scriptsize$\pm$0.028} & 0.355{\scriptsize$\pm$0.019} \\
banknote & 0.071{\scriptsize$\pm$0.009} & 0.452{\scriptsize$\pm$0.022} & 0.419{\scriptsize$\pm$0.046} & 0.453{\scriptsize$\pm$0.041} & 0.445{\scriptsize$\pm$0.038} & \textbf{0.413{\scriptsize$\pm$0.030}} \\
qsar & 0.474{\scriptsize$\pm$0.076} & 0.595{\scriptsize$\pm$0.057} & 0.567{\scriptsize$\pm$0.062} & \textbf{0.528{\scriptsize$\pm$0.047}} & 0.544{\scriptsize$\pm$0.041} & 0.541{\scriptsize$\pm$0.048} \\
pima & 0.563{\scriptsize$\pm$0.058} & 0.612{\scriptsize$\pm$0.035} & 0.618{\scriptsize$\pm$0.037} & 0.611{\scriptsize$\pm$0.036} & 0.607{\scriptsize$\pm$0.011} & \textbf{0.606{\scriptsize$\pm$0.042}} \\
\bottomrule
\end{tabular}
}
\end{table*}

\begin{table*}[t]
  \centering
  \caption{Accuracy on the incomplete test points, 10\,\% missing.}
  \label{tab:real-accuracy-10}
  \resizebox{0.9\textwidth}{!}{\begin{tabular}{lcccccc}
\toprule
Dataset & No missing & Mean imp. & Cond.\ mean imp. & MICE & Generative & PCLR \\
\midrule
magic & 0.782{\scriptsize$\pm$0.005} & 0.763{\scriptsize$\pm$0.004} & 0.771{\scriptsize$\pm$0.002} & 0.769{\scriptsize$\pm$0.005} & 0.758{\scriptsize$\pm$0.003} & \textbf{0.771{\scriptsize$\pm$0.002}} \\
phoneme & 0.721{\scriptsize$\pm$0.009} & 0.713{\scriptsize$\pm$0.010} & 0.718{\scriptsize$\pm$0.009} & 0.715{\scriptsize$\pm$0.008} & \textbf{0.750{\scriptsize$\pm$0.006}} & 0.720{\scriptsize$\pm$0.008} \\
spambase & 0.921{\scriptsize$\pm$0.009} & 0.914{\scriptsize$\pm$0.009} & 0.912{\scriptsize$\pm$0.009} & 0.910{\scriptsize$\pm$0.010} & 0.840{\scriptsize$\pm$0.027} & \textbf{0.918{\scriptsize$\pm$0.010}} \\
banknote & 0.976{\scriptsize$\pm$0.004} & 0.932{\scriptsize$\pm$0.007} & 0.944{\scriptsize$\pm$0.009} & \textbf{0.945{\scriptsize$\pm$0.006}} & 0.863{\scriptsize$\pm$0.041} & 0.943{\scriptsize$\pm$0.010} \\
qsar & 0.814{\scriptsize$\pm$0.015} & 0.782{\scriptsize$\pm$0.020} & 0.786{\scriptsize$\pm$0.016} & \textbf{0.816{\scriptsize$\pm$0.013}} & 0.700{\scriptsize$\pm$0.027} & 0.783{\scriptsize$\pm$0.014} \\
pima & 0.734{\scriptsize$\pm$0.053} & 0.729{\scriptsize$\pm$0.041} & \textbf{0.731{\scriptsize$\pm$0.046}} & 0.725{\scriptsize$\pm$0.057} & 0.661{\scriptsize$\pm$0.021} & 0.727{\scriptsize$\pm$0.049} \\
\bottomrule
\end{tabular}
}
\end{table*}

\begin{table*}[t]
  \centering
  \caption{Accuracy on the incomplete test points, 30\,\% missing.}
  \label{tab:real-accuracy-30}
  \resizebox{0.9\textwidth}{!}{\begin{tabular}{lcccccc}
\toprule
Dataset & No missing & Mean imp. & Cond.\ mean imp. & MICE & Generative & PCLR \\
\midrule
magic & 0.782{\scriptsize$\pm$0.005} & 0.735{\scriptsize$\pm$0.004} & 0.755{\scriptsize$\pm$0.006} & 0.743{\scriptsize$\pm$0.005} & 0.751{\scriptsize$\pm$0.001} & \textbf{0.756{\scriptsize$\pm$0.004}} \\
phoneme & 0.721{\scriptsize$\pm$0.009} & 0.693{\scriptsize$\pm$0.016} & 0.705{\scriptsize$\pm$0.013} & 0.701{\scriptsize$\pm$0.012} & \textbf{0.736{\scriptsize$\pm$0.013}} & 0.709{\scriptsize$\pm$0.012} \\
spambase & 0.921{\scriptsize$\pm$0.009} & \textbf{0.894{\scriptsize$\pm$0.011}} & 0.884{\scriptsize$\pm$0.006} & 0.886{\scriptsize$\pm$0.006} & 0.845{\scriptsize$\pm$0.016} & 0.884{\scriptsize$\pm$0.010} \\
banknote & 0.976{\scriptsize$\pm$0.004} & 0.836{\scriptsize$\pm$0.015} & 0.865{\scriptsize$\pm$0.020} & 0.866{\scriptsize$\pm$0.005} & 0.844{\scriptsize$\pm$0.018} & \textbf{0.875{\scriptsize$\pm$0.016}} \\
qsar & 0.814{\scriptsize$\pm$0.015} & 0.759{\scriptsize$\pm$0.026} & 0.776{\scriptsize$\pm$0.015} & \textbf{0.791{\scriptsize$\pm$0.017}} & 0.738{\scriptsize$\pm$0.033} & 0.780{\scriptsize$\pm$0.011} \\
pima & 0.734{\scriptsize$\pm$0.053} & 0.691{\scriptsize$\pm$0.042} & 0.701{\scriptsize$\pm$0.030} & 0.701{\scriptsize$\pm$0.043} & 0.657{\scriptsize$\pm$0.029} & \textbf{0.709{\scriptsize$\pm$0.033}} \\
\bottomrule
\end{tabular}
}
\end{table*}

\begin{table*}[t]
  \centering
  \caption{Accuracy on the incomplete test points, 50\,\% missing.}
  \label{tab:real-accuracy-50}
  \resizebox{0.9\textwidth}{!}{\begin{tabular}{lcccccc}
\toprule
Dataset & No missing & Mean imp. & Cond.\ mean imp. & MICE & Generative & PCLR \\
\midrule
magic & 0.782{\scriptsize$\pm$0.005} & 0.699{\scriptsize$\pm$0.008} & 0.741{\scriptsize$\pm$0.008} & 0.715{\scriptsize$\pm$0.008} & 0.742{\scriptsize$\pm$0.002} & \textbf{0.745{\scriptsize$\pm$0.007}} \\
phoneme & 0.721{\scriptsize$\pm$0.009} & 0.654{\scriptsize$\pm$0.015} & 0.679{\scriptsize$\pm$0.015} & 0.661{\scriptsize$\pm$0.018} & \textbf{0.730{\scriptsize$\pm$0.009}} & 0.700{\scriptsize$\pm$0.011} \\
spambase & 0.921{\scriptsize$\pm$0.009} & \textbf{0.870{\scriptsize$\pm$0.014}} & 0.848{\scriptsize$\pm$0.011} & 0.869{\scriptsize$\pm$0.011} & 0.820{\scriptsize$\pm$0.017} & 0.850{\scriptsize$\pm$0.010} \\
banknote & 0.976{\scriptsize$\pm$0.004} & 0.737{\scriptsize$\pm$0.017} & \textbf{0.781{\scriptsize$\pm$0.025}} & 0.753{\scriptsize$\pm$0.023} & 0.777{\scriptsize$\pm$0.027} & 0.778{\scriptsize$\pm$0.013} \\
qsar & 0.814{\scriptsize$\pm$0.015} & 0.743{\scriptsize$\pm$0.009} & 0.762{\scriptsize$\pm$0.016} & 0.763{\scriptsize$\pm$0.023} & 0.696{\scriptsize$\pm$0.048} & \textbf{0.768{\scriptsize$\pm$0.015}} \\
pima & 0.734{\scriptsize$\pm$0.053} & 0.653{\scriptsize$\pm$0.040} & 0.664{\scriptsize$\pm$0.040} & 0.673{\scriptsize$\pm$0.024} & 0.655{\scriptsize$\pm$0.021} & \textbf{0.675{\scriptsize$\pm$0.043}} \\
\bottomrule
\end{tabular}
}
\end{table*}

\section{Related Work}
\label{sec:related-work}

\paragraph{Prediction with missing features.}
The most common way to handle missing features is to impute them, with the mean, by chained equations \citep{mice}, random forests \citep{missforest} or deep generative models \citep{yoon2018gain, mattei2019miwae, ipsen2021notmiwae}, and then apply the classifier.
A single imputation ignores the uncertainty about the missing features.
Other approaches train a separate model per pattern of missing features \citep{saartsechansky2007, stempfle2025patternbypattern} or a network that adapts to the pattern \citep{lemorvan2020neumiss}; \citet{ibrahim2005glm} review the methods for generalised linear models.
We instead marginalise the missing features under an explicit density.

\paragraph{Expected predictions.}
\citet{khosravi2019expect} show that the expected prediction \eqref{eq:goal} of an LR is NP-hard even under naive Bayes, and make it tractable by fitting a naive Bayes over binary features that is conformant with the LR, so that the expectation becomes a marginal query.
This keeps the feature distribution simple and has to be refitted whenever the classifier changes; we keep a rich density over continuous features and rewrite the sigmoid instead.
\citet{jiang2020logreg} fit a joint Gaussian model of the features together with an LR from incomplete data with a stochastic EM algorithm and predict with MC.

% TODO: Characteristic Circuits (Yu et al. 2023), Polynomial Semantics of PCs (Broadrick et al. 2024), Gil-Pelaez (1951).

\section{Conclusion}
\label{sec:conclusion}

% TODO: rewrite for CFI and PCLR.
